\subsection{Registered connected-host results}
\label{sec:alpha7results}

The connected $\alpha=7$ construction fixes the dimensionless curvature
normalization and the two curvature couplings,
\[
 \frac{G_0}{a^2}=\frac{\sqrt3}{16\ln2}=0.1561763194,\qquad
 (\kappa_0,\Delta)=(2.20205311516,0.326138328034).
\]
The associated geometry weight is sampled with the capacity gate off.
The cosmological fugacity and volume feedback remain ensemble controls.
The registration of 17 October 2026 and its amendment were frozen before
sampling at these couplings. The amendment made $N_{41}=32{,}000$ the
primary volume and retained the 16k and 64k runs as reported-only checks.
The target counts $N_{41}$ simplices; the total number of four-simplices is
$N_4=N_{41}+N_{32}$.

\paragraph{Registered observables.}
Q1 tests whether the de~Sitter-like blob is independent of its starting
configuration. Q2 compares the Newton coefficient inferred from the
homogeneous volume stiffness with the boundary coefficient $G_0$:
\[
 R\equiv\frac{G_{\rm vol}}{G_0}
 \simeq 0.076503\,\frac{\widetilde s_0^{4/3}}{k_1}.
\]
The width $\widetilde s_0$ and bias-corrected local stiffness $k_1$ are
measured; their conversion is fixed by the registered dictionary. The
ratio is independent of the chain constant $m_4$. No
renormalized stiffness or source response was tuned to make $R=1$.
The written prior was $R\simeq1.4$, with disagreement considered more likely.

The final window contains 200,000 sweeps per run. Blob eligibility requires
a relative profile-fit error at most $0.05$, width $2.5\leq w\leq K/8$,
and every registered half-window stationarity statistic to satisfy
$|z|\leq3$. Here $K=80$ is the number of time slices. Q2 uses eligible
32k runs only. Its registered agreement rule places the central geometric
mean in $0.80\leq\overline R\leq1.25$; it does not require the entire
uncertainty interval to lie inside that band.

\begin{center}
\small
\begin{tabular}{lrrrll}
\hline
Run & $N_{41}$ target & $R$ & $\sigma_{\ln R}$ & Blob eligibility & Role\\
\hline
A3 & 32,000 & 1.198 & 0.059 & Drifting & Primary pool\\
A4 & 32,000 & 1.064 & 0.089 & Pass & Primary pool\\
A6 & 32,000 & 1.027 & 0.062 & Drifting & Primary pool\\
A1 & 16,000 & 1.164 & 0.061 & Drifting & Reported only\\
A2 & 16,000 & 1.060 & 0.065 & Pass & Reported only\\
A5 & 64,000 & 0.893 & 0.054 & Pass & Reported only\\
\hline
\end{tabular}
\end{center}

\paragraph{Normalization result.}
A4 is the sole eligible primary run. It gives
\[
 \overline R=1.064,\qquad \sigma_{\ln\overline R}=0.089,
 \qquad \text{Q2: AGREE}.
\]
The approximate two-standard-error interval is $0.89$--$1.27$.
Its relative profile-fit error is $0.013$. A2 and A5 also have stationary
de~Sitter-like profiles, with fit errors $0.015$ and $0.012$, and their
ratios lie inside the same agreement band. They support the volume
comparison without changing the primary verdict. Each eligible volume
has one stationary run. The change from A4 to A5 is $1.7$ standard errors
in $\ln R$; a large-volume limit has not been estimated from these data.
The specified microscopic history measure thus has a finite-volume
stiffness compatible with its own boundary normalization under the
registered conversion.

\paragraph{Phase and estimator checks.}
Q1 remains \textsf{undetermined}. A3 fails stationarity through $\xi$
at $|z|=3.3$, and A6 through $\widetilde s_0$ at $|z|=3.3$.
A3 and A4 also differ in largest-slice fraction by $6.1$ standard errors.
The existence of stationary blobs therefore coexists with an unresolved
test of start independence. A post hoc inverse-covariance calculation,
after removing the active small-sample correction floor, gives
$R=1.179,1.022,0.973$ for A2, A4 and A5. These are exploratory estimator
checks; the local-stiffness values above determine the registered result.

\paragraph{Logged continuation and the bifurcation boundary.}
The published transition between the bifurcation and de~Sitter phases of
the bare CDT measure at $\kappa_0=2.2$ and $80$ slices moves to larger
$\Delta$ as the volume grows, with pseudo-critical values $0.22$, $0.30$,
$0.34$ and $0.355\pm0.01$ at $N_{41}=40$k, 80k, 140k and 160k
\cite{coumbe,higherorder}. The published finite-size form
$\Delta_c(N)=\Delta_c(\infty)-AN^{-1/\gamma}$ with $\gamma=2.71$, fitted to
these four values with errors $0.01$, gives
$\Delta_c(\infty)=0.552\pm0.024$ and reaches the host value
$\Delta=0.326$ at $N_{41}\simeq1.1\times10^5$. The registered volumes lie
above the fitted boundary by $0.24$, $0.13$ and $0.05$ at 16k, 32k and
64k, on the de~Sitter side. The connected patches reproduce the plain CDT
weight exactly, so the published scaling applies to this measure, and it
places the host couplings in the bifurcation phase at large volume.

A logged continuation samples the first volume beyond the fitted
crossing. Its windows and reading were fixed before the runs. It extends
A2 and A5 by $200{,}000$ sweeps each and grows the final A5 state to
$N_{41}=128{,}000$, where it runs for $300{,}000$ sweeps. All $140$ engine
audits pass. At 128k the fitted boundary lies at
$\Delta=0.337\pm0.006$, above the host value by $0.011$, a margin
comparable to the $0.01$ uncertainty of the published points.

The published order parameter of the transition compares the largest
vertex order of neighboring slices \cite{coumbe}. Let $O(t)$ be the
largest vertex order in slice $t$. The central window is the run of
consecutive slices around the slice with the most vertices whose vertex
counts are at least half of that maximum, trimmed to an even length. The
alternation is
\[
 q=\frac{\bigl|\sum_i(-1)^iO(t_i)\bigr|}{\sum_iO(t_i)},
\]
with the sums over the window. Checkpoints every $2{,}500$ sweeps, after
the first $20{,}000$ new sweeps at 16k and 64k and after the first
$50{,}000$ sweeps at the 128k target, give the following means, with
errors from ten blocks.

\begin{center}
\small
\begin{tabular}{lrrll}
\hline
Run & $N_{41}$ & Checkpoints & $\overline q$ & $\overline{O_{\max}/N_4}$\\
\hline
A2 continued & 16,000 & 72 & $0.061\pm0.008$ & $0.0493\pm0.0011$\\
A5 continued & 64,000 & 72 & $0.082\pm0.008$ & $0.0303\pm0.0006$\\
B128 & 128,000 & 100 & $0.075\pm0.006$ & $0.0205\pm0.0003$\\
\hline
\end{tabular}
\end{center}

The registered reading is no growth. The 64k mean exceeds the 16k mean by
$1.9$ combined errors, below the registered threshold of two, and the
128k mean lies below the 64k mean. The largest order grows as
$N_4^{0.59}$, slower than the volume, so no vertex condenses. The
registered bifurcation criterion, $\overline q\geq0.15$ at 128k with the
largest order growing at least as $N_4^{0.90}$, fails on both counts. In
the 64k and 128k runs the sign of the alternation never changes, and both
runs descend from one 64k state, so their block errors are lower bounds.
The 128k mean exceeds the 16k mean by $1.4$ combined errors, and the 16k
point lies $0.24$ above the fitted boundary.

The last $200{,}000$ sweeps of each continuation run give reported-only
ratios.

\begin{center}
\small
\begin{tabular}{lrrrl}
\hline
Run & $N_{41}$ & $R$ & $\sigma_{\ln R}$ & Blob eligibility\\
\hline
A2 continued & 16,000 & 0.858 & 0.059 & Pass\\
A5 continued & 64,000 & 1.001 & 0.061 & Pass\\
B128 & 128,000 & 1.045 & 0.045 & Drifting, width $11.9$\\
\hline
\end{tabular}
\end{center}

All three ratios lie inside the agreement band. Successive windows of one
run differ by $0.21$ in $\ln R$ at 16k and by $0.11$ at 64k, more than the
jackknife error of a single window, so the scatter between windows sets
the uncertainty of a reported-only value. The 128k window fails
stationarity in $\widetilde s_0$ at $|z|=3.8$, and its width exceeds the
registered maximum $K/8=10$. No reading was fixed in advance for these
ratios.

The phase at larger volume is open. The fitted boundary exceeds the host
value by $0.04$ at $N_{41}=2\times10^5$ and by $0.06$ at $2.6\times10^5$,
and runs at those volumes test whether the boundary of this measure
follows the published form.

\paragraph{Matter-response status.}
A separately registered curvature-trace test used two new streams at
fixed $\kappa_4$, with 64 saved geometries and two source marks per
geometry in each stream. Its dimensionless diagnostic was
$\gamma_{\rm trace}=1.82\pm1.04$, with approximate 95\% interval
$[-0.65,4.28]$, against the target $G_0/a^2=0.1562$.
Both streams also exceeded the registered maximum blob width of 10,
with widths $11.17$ and $10.61$. The verdict is
\textsf{undetermined}; no matter Newton coefficient was reported.
This separate fixed-fugacity gate failure leaves Q2 unchanged.

Subsequent quenched-scalar measurements on those 128 saved geometries
are development analyses. They use both saved sources and six bare
masses. The inferred binding slopes depend on distance range, binning
and scalar discretization, while fixed nongravitating controls also
produce apparent short-range slopes. A conditional-variance diagnostic
does not restore a stable binding interpretation. These measurements
supply no accepted binding energy or gravitational null result.
Further matter-response production is deferred pending a calibrated
finite-geometry observable and an independent power assessment.

The completed evidence concerns the microscopic realization, stationary
finite-volume profiles and the registered homogeneous stiffness ratio.
Equality with a matter-source Newton coefficient, a common matter cone,
and the interacting massless transverse--traceless sector remain open.
The bare ratio $G_0/G_*=0.9975\pm0.0010$ rests on the face-wise lock and the event calibration. With $R$ it gives $G_{\rm vol}/G_*=1.061$, with logarithmic standard error $0.089$.
